\section{Proof of ReLU Dense Span Property for Frozen Random Features}
\label{sec:relu-dense-span-proof}

This section provides a rigorous proof that frozen random features with full support give dense span in $L^2(\mathcal{P}_X)$, adapting the approach from the full-width case \cite{nguyen2023rigorousframeworkmeanfield}.

\subsection{Main Result}

\begin{theorem}[Dense span for frozen ReLU random features]
\label{thm:relu-dense-span}
Let $\varphi_1$ be the ReLU activation function, and let $f_1(C_1)$ be frozen random feature vectors with $\text{supp}(f_1(C_1)) = \mathbb{R}^d$. Then the function class:
\[
\mathcal{F} = \{\varphi_1(\langle f_1(c_1), \cdot\rangle) : c_1 \in \Omega_1\}
\]
has dense span in $L^2(\mathcal{P}_X)$ (with respect to the data distribution $\mathcal{P}_X$).

In other words, for any $g \in L^2(\mathcal{P}_X)$ and any $\epsilon > 0$, there exist coefficients $\alpha_1, \ldots, \alpha_n$ and feature vectors $c_1, \ldots, c_n \in \Omega_1$ such that:
\[
\left\|g - \sum_{i=1}^n \alpha_i \varphi_1(\langle f_1(c_i), \cdot\rangle)\right\|_{L^2(\mathcal{P}_X)} < \epsilon.
\]
\end{theorem}

\subsection{Key Insight: Simpler Than Full-Width Case}

\textbf{Advantage of frozen features:} Unlike the full-width case where we need to prove that full support is \emph{maintained} during training (using algebraic topology, Lemma~\ref{lem:full-support-3} in \cite{nguyen2023rigorousframeworkmeanfield}), in our case:
\begin{itemize}
\item $f_1(C_1)$ are \textbf{frozen} (not trained), so $\text{supp}(f_1(C_1)) = \mathbb{R}^d$ trivially
\item We only need to prove that full support \textbf{implies} dense span
\item This is a standard universal approximation result, not a dynamical property
\end{itemize}

\subsection{Proof Strategy}

The proof proceeds in three steps:

\begin{enumerate}
\item \textbf{Standard universal approximation result:} Use Leshno et al. (1993) to show that $\{\varphi_1(\langle u, \cdot\rangle) : u \in \mathbb{R}^d\}$ has dense span in $L^2(\mathcal{P}_X)$.

\item \textbf{Full support gives dense parameterization:} Since $\text{supp}(f_1(C_1)) = \mathbb{R}^d$, the set $\{f_1(c_1) : c_1 \in \Omega_1\}$ is dense in $\mathbb{R}^d$.

\item \textbf{Continuity argument:} Use continuity of $u \mapsto \varphi_1(\langle u, \cdot\rangle)$ in $L^2(\mathcal{P}_X)$ to show that dense parameterization implies dense span.
\end{enumerate}

\subsection{Step 1: Standard Universal Approximation}

\begin{lemma}[Leshno et al. 1993, adapted]
\label{lem:leshno-relu}
Let $\varphi_1$ be ReLU (a non-polynomial activation). Then the function class:
\[
\mathcal{G} = \{\varphi_1(\langle u, \cdot\rangle) : u \in \mathbb{R}^d\}
\]
has dense span in $L^2(\mathcal{P}_X)$.

More precisely, for any $g \in L^2(\mathcal{P}_X)$ and $\epsilon > 0$, there exist $u_1, \ldots, u_n \in \mathbb{R}^d$ and coefficients $\alpha_1, \ldots, \alpha_n \in \mathbb{R}$ such that:
\[
\left\|g - \sum_{i=1}^n \alpha_i \varphi_1(\langle u_i, \cdot\rangle)\right\|_{L^2(\mathcal{P}_X)} < \epsilon.
\]
\end{lemma}

\begin{proof}[Proof sketch]
This follows from the universal approximation theorem of Leshno et al. (1993), which states that single-hidden-layer networks with non-polynomial activations are dense in $C(\mathbb{R}^d)$ (and hence in $L^2$ for compactly supported measures).

For ReLU specifically:
\begin{itemize}
\item ReLU is non-polynomial (it's piecewise linear but not a polynomial)
\item The result extends to $L^2(\mathcal{P}_X)$ when $\mathcal{P}_X$ has compact support (which we assume: $|X| \le K$ with probability 1)
\item The linear combinations $\sum_{i=1}^n \alpha_i \varphi_1(\langle u_i, \cdot\rangle)$ can approximate any function in $L^2(\mathcal{P}_X)$
\end{itemize}

For a complete proof, see Leshno et al. (1993, Theorem 1) or Cybenko (1989) adapted to ReLU.
\end{proof}

\subsection{Step 2: Full Support Implies Dense Parameterization}

\begin{lemma}[Full support gives dense parameterization]
\label{lem:dense-parameterization}
If $\text{supp}(f_1(C_1)) = \mathbb{R}^d$, then the set $\{f_1(c_1) : c_1 \in \Omega_1\}$ is dense in $\mathbb{R}^d$.

In other words, for any $u \in \mathbb{R}^d$ and any $\delta > 0$, there exists $c_1 \in \Omega_1$ such that $|f_1(c_1) - u| < \delta$.
\end{lemma}

\begin{proof}
This is immediate from the definition of support. If $\text{supp}(f_1(C_1)) = \mathbb{R}^d$, then for any open ball $B_\delta(u) \subset \mathbb{R}^d$, we have:
\[
\mathbb{P}(f_1(C_1) \in B_\delta(u)) > 0.
\]
This means there exists at least one $c_1 \in \Omega_1$ such that $f_1(c_1) \in B_\delta(u)$, i.e., $|f_1(c_1) - u| < \delta$.

Since this holds for any $u \in \mathbb{R}^d$ and any $\delta > 0$, the set $\{f_1(c_1) : c_1 \in \Omega_1\}$ is dense in $\mathbb{R}^d$.
\end{proof}

\subsection{Step 3: Continuity Argument}

\begin{lemma}[Continuity of ReLU features]
\label{lem:continuity-relu}
The mapping $u \mapsto \varphi_1(\langle u, \cdot\rangle)$ is continuous from $\mathbb{R}^d$ to $L^2(\mathcal{P}_X)$.

More precisely, for any $u, u' \in \mathbb{R}^d$:
\[
\|\varphi_1(\langle u, \cdot\rangle) - \varphi_1(\langle u', \cdot\rangle)\|_{L^2(\mathcal{P}_X)} \le K |u - u'|,
\]
where $K$ depends on the bound on $|X|$ and the Lipschitz constant of $\varphi_1$.
\end{lemma}

\begin{proof}
Since $\varphi_1$ is ReLU (which is 1-Lipschitz) and $|X| \le K$ with probability 1, we have:
\begin{align*}
\|\varphi_1(\langle u, \cdot\rangle) - \varphi_1(\langle u', \cdot\rangle)\|_{L^2(\mathcal{P}_X)}^2 
&= \mathbb{E}_X[|\varphi_1(\langle u, X\rangle) - \varphi_1(\langle u', X\rangle)|^2]\\
&\le \mathbb{E}_X[|\langle u - u', X\rangle|^2] \quad \text{(ReLU is 1-Lipschitz)}\\
&\le |u - u'|^2 \mathbb{E}_X[|X|^2]\\
&\le K^2 |u - u'|^2.
\end{align*}
Taking square roots gives the desired bound.
\end{proof}

\subsection{Completing the Proof of Theorem~\ref{thm:relu-dense-span}}

\begin{proof}[Proof of Theorem~\ref{thm:relu-dense-span}]
By Lemma~\ref{lem:leshno-relu}, for any $g \in L^2(\mathcal{P}_X)$ and $\epsilon > 0$, there exist $u_1, \ldots, u_n \in \mathbb{R}^d$ and coefficients $\alpha_1, \ldots, \alpha_n$ such that:
\[
\left\|g - \sum_{i=1}^n \alpha_i \varphi_1(\langle u_i, \cdot\rangle)\right\|_{L^2(\mathcal{P}_X)} < \frac{\epsilon}{2}.
\]

By Lemma~\ref{lem:dense-parameterization}, since $\text{supp}(f_1(C_1)) = \mathbb{R}^d$, for each $u_i$ and $\delta > 0$ (to be chosen), there exists $c_i \in \Omega_1$ such that $|f_1(c_i) - u_i| < \delta$.

By Lemma~\ref{lem:continuity-relu}, we have:
\[
\|\varphi_1(\langle f_1(c_i), \cdot\rangle) - \varphi_1(\langle u_i, \cdot\rangle)\|_{L^2(\mathcal{P}_X)} \le K |f_1(c_i) - u_i| < K\delta.
\]

Therefore:
\begin{align*}
\left\|g - \sum_{i=1}^n \alpha_i \varphi_1(\langle f_1(c_i), \cdot\rangle)\right\|_{L^2(\mathcal{P}_X)}
&\le \left\|g - \sum_{i=1}^n \alpha_i \varphi_1(\langle u_i, \cdot\rangle)\right\|_{L^2(\mathcal{P}_X)}\\
&\quad + \left\|\sum_{i=1}^n \alpha_i (\varphi_1(\langle u_i, \cdot\rangle) - \varphi_1(\langle f_1(c_i), \cdot\rangle))\right\|_{L^2(\mathcal{P}_X)}\\
&< \frac{\epsilon}{2} + \sum_{i=1}^n |\alpha_i| K\delta.
\end{align*}

Choosing $\delta = \frac{\epsilon}{2K \sum_{i=1}^n |\alpha_i|}$ (which is finite since we can choose the $\alpha_i$ to be bounded), we get:
\[
\left\|g - \sum_{i=1}^n \alpha_i \varphi_1(\langle f_1(c_i), \cdot\rangle)\right\|_{L^2(\mathcal{P}_X)} < \epsilon.
\]

This proves that $\mathcal{F}$ has dense span in $L^2(\mathcal{P}_X)$.
\end{proof}

\subsection{Connection to Full-Width Proof}

The full-width proof (lines 2725-2726 of \cite{nguyen2023rigorousframeworkmeanfield}) uses the same dense span property, but they need to prove that full support is \emph{maintained} during training using an algebraic topology argument (Lemma~\ref{lem:full-support-3}).

In our case:
\begin{itemize}
\item \textbf{No algebraic topology needed:} Since $f_1(C_1)$ are frozen, full support is automatic
\item \textbf{Direct application:} We can directly apply Leshno et al. (1993) + continuity argument
\item \textbf{Simpler proof:} The proof is more straightforward because we don't need to track support evolution
\end{itemize}

This is a \textbf{key theoretical advantage} of the frozen random feature approach: it eliminates the need for complex support-preservation arguments while still maintaining universal approximation.

\subsection{Remarks}

\begin{enumerate}
\item \textbf{Non-polynomial requirement:} The proof relies on $\varphi_1$ being non-polynomial. ReLU satisfies this (it's piecewise linear but not a polynomial).

\item \textbf{Data distribution:} The result holds for any $\mathcal{P}_X$ with bounded support ($|X| \le K$). For unbounded support, additional conditions may be needed.

\item \textbf{Infinite-width limit:} The proof works in the infinite-width limit where $\Omega_1$ is a continuous space, as long as $\text{supp}(f_1(C_1)) = \mathbb{R}^d$.

\item \textbf{Extension to other activations:} The same argument works for any non-polynomial activation (sigmoid, tanh, etc.), not just ReLU.
\end{enumerate}
